%! Modeling with Quadratics — Unit 2, Lesson 2.7 — Slides (Beamer)
\documentclass[aspectratio=169, 11pt]{beamer}
\usepackage{algebra2-beamer}   % provides \forestheader{} and \sectionlabel[color]{}

\begin{document}

% TITLE SLIDE (CourseName is not defined in beamer, so the name is literal)
{
\setbeamercolor{background canvas}{bg=forest}
\begin{frame}[plain]
  \vspace{2.8cm}\hspace{0.55cm}%
  \begin{minipage}{0.9\textwidth}
    {\color{goldacc}\bfseries\small LESSON 2.7}\\[0.5cm]
    {\color{white}\bfseries\LARGE Modeling with Quadratics}\\[1.2cm]
    {\color{forestmid}\small Algebra 2: Shepherd~$\cdot$~Unit 2: Quadratic Functions}
  \end{minipage}
\end{frame}
}

% HOOK — three questions, one parabola
\begin{frame}[t, plain]
  \forestheader{Every question is really about a feature}
  \sectionlabel{Hook}
  \begin{columns}[T]
    \begin{column}{0.54\textwidth}
      A ball thrown up at $32$ ft/s from a $48$-ft ledge:
      \[ h(t)=-16t^2+32t+48. \]
      \begin{itemize}\setlength{\itemsep}{4pt}
        \item ``How high does it \textbf{start}?'' $\to$ $y$-intercept
        \item ``\textbf{Highest} point, and when?'' $\to$ vertex
        \item ``When does it \textbf{land}?'' $\to$ a zero
      \end{itemize}
    \end{column}
    \begin{column}{0.46\textwidth}
      \centering
      \begin{tikzpicture}[xscale=1.15, yscale=0.045]
        \draw[->, thick, charcoal] (-0.15,0)--(3.6,0) node[below right, font=\tiny]{$t$};
        \draw[->, thick, charcoal] (0,-3)--(0,74) node[left, font=\tiny]{$h$};
        \draw[very thick, forest, smooth, samples=60, domain=0:3] plot (\x,{-16*\x*\x+32*\x+48});
        \draw[dashed, linegray] (1,0)--(1,64)--(0,64);
        \node[circle, fill=charcoal, inner sep=1.3pt] at (0,48){};
        \node[circle, fill=charcoal, inner sep=1.3pt] at (1,64){};
        \node[circle, fill=charcoal, inner sep=1.3pt] at (3,0){};
        \node[font=\tiny, anchor=west] at (1.06,64){$(1,64)$};
      \end{tikzpicture}
    \end{column}
  \end{columns}
\end{frame}

% PROJECTILE — feature map
\begin{frame}[t, plain]
  \forestheader{Projectile: $h(t)=-16t^2+v_0t+h_0$}
  \sectionlabel[goldacc]{Model 1}
  \begin{columns}[T]
    \begin{column}{0.5\textwidth}
      \begin{block}{The feature map}
        \textbf{Start} $=h(0)$ ($y$-int) \\
        \textbf{Max/when} $=$ vertex ($t=-\tfrac{b}{2a}$, then evaluate) \\
        \textbf{Lands} $=$ positive zero ($h=0$)
      \end{block}
    \end{column}
    \begin{column}{0.5\textwidth}
      For $h(t)=-16t^2+32t+48$:
      \begin{itemize}\setlength{\itemsep}{2pt}
        \item Start: $h(0)=48$ ft
        \item Vertex: $t=1$, $h(1)=64$ ft
        \item Lands: $t^2-2t-3=0\Rightarrow t=3$ (reject $-1$)
      \end{itemize}
      \small Input $=$ \emph{when}; output $=$ \emph{how high}.
    \end{column}
  \end{columns}
\end{frame}

% MAX AREA
\begin{frame}[t, plain]
  \forestheader{Maximum area: the vertex is the ``most''}
  \sectionlabel{Model 2}
  \begin{columns}[T]
    \begin{column}{0.54\textwidth}
      Pen against a barn wall, $40$ ft of fence, \emph{three} sides. Width $x$, length $40-2x$:
      \[ A(x)=x(40-2x)=-2x^2+40x. \]
      \[ x=-\tfrac{b}{2a}=10\ \Rightarrow\ A=200\text{ ft}^2. \]
      Dimensions $10\times20$; feasible domain $0<x<20$.
    \end{column}
    \begin{column}{0.46\textwidth}
      \centering
      \begin{tikzpicture}[scale=0.8]
        \draw[very thick, charcoal] (-0.3,1.5)--(4.3,1.5);
        \node[font=\tiny, charcoal, anchor=south] at (2,1.55){barn wall};
        \draw[very thick, forest] (0,0)--(0,1.5);
        \draw[very thick, forest] (4,0)--(4,1.5);
        \draw[very thick, forest] (0,0)--(4,0);
        \node[font=\tiny, forest, anchor=east] at (-0.05,0.75){$x$};
        \node[font=\tiny, forest, anchor=west] at (4.05,0.75){$x$};
        \node[font=\tiny, forest, anchor=north] at (2,-0.05){$40-2x$};
      \end{tikzpicture}
    \end{column}
  \end{columns}
\end{frame}

% MAX REVENUE
\begin{frame}[t, plain]
  \forestheader{Optimize revenue: the best price}
  \sectionlabel[goldacc]{Model 3}
  \begin{columns}[T]
    \begin{column}{0.5\textwidth}
      $200$ tickets at \$8; each \$1 increase loses $20$ buyers ($x=$ \$1 increases).
      \[ R(x)=(8+x)(200-20x) \]
      \[ =-20x^2+40x+1600. \]
    \end{column}
    \begin{column}{0.5\textwidth}
      \begin{block}{Best price $=$ vertex}
        $x=-\tfrac{40}{2(-20)}=1$ \\[2pt]
        Charge \$9, sell $180$, max revenue \textbf{\$1620}.
      \end{block}
      \small ``Best/most/least'' always lives at the vertex.
    \end{column}
  \end{columns}
\end{frame}

% STRATEGY
\begin{frame}[t, plain]
  \forestheader{The strategy --- every time}
  \sectionlabel{Method}
  \begin{columns}[T]
    \begin{column}{0.5\textwidth}
      \begin{enumerate}\setlength{\itemsep}{3pt}
        \item Define the variable (units).
        \item Write the quadratic model.
        \item Which feature does the question want?
        \item Find it (Lessons 2.1--2.5).
        \item Interpret with units; reject impossible values.
      \end{enumerate}
    \end{column}
    \begin{column}{0.5\textwidth}
      \begin{block}{The feature-to-question map}
        \textbf{start} $\to$ $y$-intercept \\
        \textbf{``reaches $0$''} $\to$ a zero \\
        \textbf{``the most / the best''} $\to$ the vertex
      \end{block}
    \end{column}
  \end{columns}
\end{frame}

% SUMMARY + NEXT
\begin{frame}[t, plain]
  \forestheader{Summary \& what's next}
  \sectionlabel{Summary}
  \begin{columns}[T]
    \begin{column}{0.55\textwidth}
      \begin{block}{Today}
        \emph{Build} the quadratic, then read the answer off a feature: $y$-intercept $=$ start, zero $=$ reaches
        $0$, vertex $=$ max/min (input $=$ when, output $=$ how much). Always interpret with units.
      \end{block}
    \end{column}
    \begin{column}{0.45\textwidth}
      \begin{block}{Next}
        \textbf{Unit 2 Review \& Test}
        \\[2pt]
        Graphing, factoring, the quadratic formula, complex numbers, systems, and modeling --- all together.
      \end{block}
    \end{column}
  \end{columns}
\end{frame}

\end{document}
